\documentclass[10pt,letterpaper]{article}

% ATS-friendly, plain-text-first resume. No icons, graphics, tables for core content,
% or template-specific assets from previous versions.
\usepackage[letterpaper,margin=0.62in]{geometry}
\usepackage[T1]{fontenc}
\usepackage[utf8]{inputenc}
\usepackage{enumitem}
\usepackage[hidelinks]{hyperref}
\usepackage{titlesec}
\usepackage{array}

\pagestyle{empty}
\setlength{\parindent}{0pt}
\setlength{\parskip}{0pt}
\setlist[itemize]{leftmargin=1.2em,itemsep=2pt,topsep=2pt,parsep=0pt,partopsep=0pt}
\titleformat{\section}{\large\bfseries}{}{0pt}{}[\vspace{-2pt}\titlerule]
\titlespacing*{\section}{0pt}{8pt}{4pt}

\newcommand{\role}[4]{%
  \textbf{#1} \hfill #2\\
  \textit{#3} \hfill \textit{#4}\vspace{1pt}
}

\newcommand{\project}[2]{\textbf{#1}: #2}

\begin{document}

{\LARGE\textbf{Adam Petrovic}}\\
Sydney, Australia \;|\; +61 400 077 049 \;|\; \href{mailto:adam@petrovic.com.au}{adam@petrovic.com.au}\\
\href{https://linkedin.com/in/adampetrovic}{linkedin.com/in/adampetrovic} \;|\; \href{https://github.com/adampetrovic}{github.com/adampetrovic} \;|\; \href{https://adampetrovic.me}{adampetrovic.me}

\section*{Senior Software Engineer -- Cloud Infrastructure, Distributed Systems, Reliability}
Senior engineer with 14+ years of experience building high-scale cloud platforms and 10+ years at Atlassian across Jira Cloud, Observability, and Transaction Data Platform. I specialise in distributed systems, storage, Kubernetes, observability, incident response, and multi-cloud infrastructure. My strongest work is turning ambiguous platform problems into reliable production systems, clear technical direction, and reusable engineering patterns adopted beyond my immediate team.

\section*{Core Strengths}
\begin{itemize}
  \item Distributed systems and storage platforms: object storage, sharding, replication, data integrity, backup, restore, cost attribution, capacity testing.
  \item Cloud and Kubernetes infrastructure: AWS, GCP, Kubernetes, KITT, Spinnaker, service descriptors, deployment automation, multi-region and multi-cloud migration.
  \item Reliability leadership: incident response, SLOs, post-incident analysis, operational readiness, on-call improvements, synthetic monitoring, observability platforms.
  \item Technical direction: RFCs, architecture trade-offs, project planning, risk management, cross-team alignment, mentoring, engineering management.
\end{itemize}

\section*{Experience}

\role{Atlassian}{Sydney, Australia}{Senior Engineer, Transaction Data Platform / Object Store}{August 2023 -- Present}
\begin{itemize}
  \item Tech-led Object Store's GCP staging migration, establishing foundational multi-cloud capability for a critical internal storage platform and delivering a functionally complete deployment four weeks ahead of the target milestone.
  \item Created the project plan, milestone model, risk tracking, and dependency-management approach for the GCP migration; the format was adopted by multiple Transaction Data Platform teams as their template for GCP readiness work.
  \item Designed and implemented production storage capabilities in Kotlin across S3-compatible object storage, data residency, encryption, backup, checksum handling, shard provisioning, and GCS integration.
  \item Authored strategic RFCs for GCS checksum handling, TCS dependency removal, seekable compression, and related Object Store architecture decisions, balancing reliability, operational simplicity, migration risk, and hosting cost.
  \item Built S3 Inventory and DynamoDB export ingestion into Socrates, enabling service-level cost attribution, object-level verification, and self-service analytics over multi-terabyte datasets.
  \item Responded to critical Object Store incidents, including implementing a Socrates pipeline during a Sev 2, proving no data loss through analysis, and converting the response into a continuous object-verification initiative for one of the platform's highest risks.
  \item Implemented point-in-time restore support for overwritten objects and investigated complex customer-impacting behaviours including 409 conflicts, retry semantics, credential-provider issues, and multi-shard rate-limiting correctness.
  \item Created \texttt{sd-diff-pipe}, a Bitbucket Pipeline that detects risky service descriptor and archetype changes before deployment; adopted by Object Store, Media, and ERS with broader team interest.
  \item Established upload/download performance SLOs grounded in production traffic patterns and built Locust-based capacity tests that scaled to 1.3M requests per minute.
  \item In FY26H1, authored 162 pull requests, reviewed 681 pull requests, created 22 Confluence pages, and opened 85 Jira tickets while continuing project leadership and incident-response work.
\end{itemize}

\role{Atlassian}{Sydney, Australia}{Engineering Manager, Observability}{October 2020 -- August 2023}
\begin{itemize}
  \item Led the Better Metrics program migrating Atlassian engineering from Datadog to SignalFx, delivering more than \$2M in annual savings while maintaining reliability of company-wide metrics workflows.
  \item Managed an Observability engineering team through vendor evaluation, migration planning, stakeholder communication, operational risk management, hiring, coaching, and performance development.
  \item Directed multi-vendor ingestion work that allowed concurrent platform evaluation while the production metrics pipeline handled millions of messages per second.
  \item Drove config-as-code migration automation to increase migration throughput, reduce manual coordination, and meet contractual migration deadlines across a large internal customer base.
  \item Kept technical involvement in architecture, incident reviews, roadmap decisions, and operational readiness while leading people-management responsibilities.
\end{itemize}

\role{Atlassian}{Sydney, Australia}{Senior Engineer, Observability}{June 2017 -- October 2020}
\begin{itemize}
  \item Built and operated Atlassian's company-wide metrics infrastructure, scaling ingestion to 180M metrics per minute for critical production monitoring and alerting.
  \item Developed Pollinator, a Go-based synthetic monitoring and post-deployment verification platform used by production services across regions to detect reliability and performance regressions.
  \item Extended GoStatsD with Prometheus-style histogram support, enabling accurate percentile and SLO measurement for hundreds of internal services.
  \item Migrated observability components toward Kafka and Apache Flink stream processing, improving scalability and reducing operational constraints in high-volume telemetry pipelines.
\end{itemize}

\role{Atlassian}{Sydney, Australia}{Senior Site Reliability Engineer, Jira Cloud}{March 2015 -- June 2017}
\begin{itemize}
  \item Led a cross-continental SRE group delivering Jira Cloud European infrastructure ahead of a public product milestone, supporting international expansion and enterprise readiness.
  \item Refactored Jira Cloud site import processes, reducing failed customer imports, on-call incidents, and recurring operational toil.
  \item Introduced and facilitated incident-response training based on Google's Wheel of Misfortune model, helping teams practise production failure scenarios before real incidents.
  \item Operated multi-tenant Jira Cloud infrastructure through a period of rapid customer, region, and reliability growth.
\end{itemize}

\role{Freelancer.com}{Sydney, Australia}{Software Engineer}{August 2011 -- February 2015}
\begin{itemize}
  \item Built and operated authentication, messaging, notification, email, analytics, and fraud-detection platform services for a global marketplace.
  \item Helped move critical systems from PHP monolith patterns toward Python and Go services using Apache Thrift, Redis, RabbitMQ, and contract-first APIs.
  \item Improved request-path latency and operational isolation by moving asynchronous work onto event-driven infrastructure.
\end{itemize}

\section*{Selected Public Work}
\begin{itemize}
  \item \project{home-ops}{Bare-metal Kubernetes homelab using Talos Linux, Flux, Helm, Kustomize, SOPS, Rook-Ceph, and VolSync. Demonstrates production-style GitOps, storage, backup, and platform operations. \href{https://github.com/adampetrovic/home-ops}{github.com/adampetrovic/home-ops}}
  \item \project{pi-telegram-bot}{TypeScript Telegram bot that orchestrates coding-agent sessions over RPC, reflecting current AI-assisted engineering workflows. \href{https://github.com/adampetrovic/pi-telegram-bot}{github.com/adampetrovic/pi-telegram-bot}}
  \item \project{dotfiles}{Chezmoi-managed development environment and automation, including TypeScript tooling for repeatable local setup. \href{https://github.com/adampetrovic/dotfiles}{github.com/adampetrovic/dotfiles}}
  \item \project{abb-cvc-extract}{Python tool extracting time-series data from Aussie Broadband CVC graph images. \href{https://github.com/adampetrovic/abb-cvc-extract}{github.com/adampetrovic/abb-cvc-extract}}
\end{itemize}

\section*{Technical Skills}
\begin{itemize}
  \item \textbf{Languages:} Kotlin, Go, Python, Java, TypeScript, JavaScript, SQL, Bash
  \item \textbf{Cloud and Infrastructure:} AWS, GCP, Kubernetes, Docker, Talos Linux, Flux, Helm, Kustomize, Spinnaker, CI/CD
  \item \textbf{Data and Storage:} S3, GCS, DynamoDB, PostgreSQL, MySQL, Redis, Kafka, Apache Flink, RabbitMQ, object storage, sharding
  \item \textbf{Reliability and Observability:} SignalFx, Datadog, Prometheus, Grafana, Splunk, StatsD, Sentry, Locust, incident response, SLOs
\end{itemize}

\section*{Education}
\textbf{The University of Sydney} \hfill Sydney, Australia\\
Bachelor of Information Technology, Computer Science \hfill March 2008 -- July 2012\\
Dean's List; distinction average.

\end{document}
